% TOP_PROBLEM: {"id": 2584, "title": "Kourovka Notebook Problem 21.75", "category_id": 17, "category": {"id": 17, "name": "group_theory", "display_name": "Group Theory", "description": "Problems about groups, group actions, representations, and related algebraic structures.", "slug": "group-theory", "order_index": 17, "created_at": "2026-07-31T15:26:25.671Z"}, "status": "open"}
% FIRST_POSTED: null
% ENTRY_KIND: statement_only
% Imported from: batch09.tex; UnsolvedMath ID 2584
\documentclass[11pt]{article}
\usepackage[margin=25mm]{geometry}
\usepackage{fontspec}\setmainfont{Latin Modern Roman}
\usepackage{amsmath,amssymb,mathtools}
\usepackage{unicode-math}\setmathfont{Latin Modern Math}
\usepackage{xurl,hyperref}
\hypersetup{colorlinks=true,urlcolor=blue}
\setcounter{secnumdepth}{0}\setlength{\parindent}{0pt}\setlength{\parskip}{5pt}\setlength{\emergencystretch}{3em}
\newcommand{\sref}[2]{\hyperlink{source:#1:#2}{[S#2]}}
\newcommand{\cataloguescope}[1]{\paragraph{Recorded target.}#1}
\begin{document}
% UnsolvedMath ID 2584; source catalogue code KOU-21.75
\setcounter{section}{2583}
\section{Kourovka Notebook Problem 21.75}\label{2584}
{\small\sffamily UnsolvedMath ID 2584 · Group Theory}
\subsection{Definitions and mathematical statement}

\paragraph{Shared notation.}$\mathbb N=\{1,2,\ldots\}$, and $\mathbb Z,\mathbb Q,\mathbb R,\mathbb C$ denote the integers, rationals, reals and complex numbers. Any other conventions are specified locally.


For $0\leq r<m$, let $r(m)=r+m\mathbb Z$. If $r_1(m_1)$ and $r_2(m_2)$ are disjoint, their class transposition sends $r_1+km_1$ to $r_2+km_2$ and vice versa for each $k\in\mathbb Z$, fixing all other integers. Let $\operatorname{CT}(\mathbb Z)$ be generated by all these permutations. For a set $P$ of odd primes, $\operatorname{CT}_P(\mathbb Z)$ is generated by the class transpositions for which each prime factor of both moduli belongs to $P\cup\{2\}$. The notation $\langle A,B\rangle$ denotes the subgroup generated jointly by two subgroups. The prime sets may be infinite.
\paragraph{Statement recorded in the supplied source.}
For every two sets $P_1,P_2$ of odd primes with $P_1\nsubseteq P_2$ and $P_2\nsubseteq P_1$, is the natural inclusion always strict:
\[
 \langle\operatorname{CT}_{P_1}(\mathbb Z),\operatorname{CT}_{P_2}(\mathbb Z)\rangle
 \lneq\operatorname{CT}_{P_1\cup P_2}(\mathbb Z)?
\] \sref{2584}{2}
\subsection{Short English statement}
Do two incomparable prime sets always generate a proper subgroup of the class-transposition group permitted by their union?
\subsection{Sources}
\begin{description}
\item[\hypertarget{source:2584:1}{S1}] ULAM.AI and the UnsolvedMath Contributors, \emph{UnsolvedMath: A Curated Collection of Open Mathematics Problems}, record 2584 (KOU-21.75). CC BY 4.0. \url{https://huggingface.co/datasets/ulamai/UnsolvedMath}.
\item[\hypertarget{source:2584:2}{S2}] S. Kohl, Problem 21.75, in E. I. Khukhro and V. D. Mazurov (eds.), \emph{The Kourovka Notebook: Unsolved Problems in Group Theory}, 21st edition, arXiv:1401.0300v48 (6 October 2026). See also source definitions in Problems 17.57 and17.60. This version records equality instead of the proposed strict containment; the displayed question is the historically registered target. \url{https://arxiv.org/abs/1401.0300v48}.
\end{description}
\label{end:2584}
\end{document}
